\clearpage
\VerdierPhysicalPageAnchor{0185}%
\begingroup\small\setlength{\parskip}{0pt}\setlength{\emergencystretch}{1.5em}
\setlength{\abovedisplayskip}{4pt plus 2pt minus 2pt}%
\setlength{\belowdisplayskip}{4pt plus 2pt minus 2pt}%

\VerdierSemanticAnchor{ch3d-1-3-6-conditions-cont}
\noindent\textit{i) $F$ se factorise par $K(\mathcal A)$.}

\noindent\textit{ii) $F$ transforme les homotopismes en isomorphismes de
$C$.}

\VerdierSemanticAnchor{ch3d-1-3-6-proof}
Il est clair que i) $\Rightarrow$ ii). Pour montrer que ii)
$\Rightarrow$ i), il suffit de montrer que $F$ transforme les couples de
morphismes homotopes en morphismes égaux. Soit $X$ un objet de
$\operatorname{comp}(\mathcal A)$. Le foncteur sur
$\operatorname{comp}(\mathcal A)$ :
\[
Y\longmapsto
\{\text{couple de morphismes de $X$ dans $Y$ muni d'une homotopie qui les
relie}\}
\]
est représentable par un complexe $Z$ :

\VerdierSemanticAnchor{ch3d-1-3-6-representing-complex}
\begin{equation}
\begin{aligned}
Z^i&=X^i\oplus X^i\oplus X^{i+1},\\[2pt]
d_Z^i&=
\begin{pmatrix}
d_X^i&0&\operatorname{id}_{X^{i+1}}\\
0&d_X^i&-\operatorname{id}_{X^{i+1}}\\
0&0&-d_X^{i+1}
\end{pmatrix}.
\end{aligned}
\tag{1.3.6.1}
\end{equation}
Il existe donc deux morphismes homotopes
$p_1,p_2\colon X\longrightarrow Z$ :

\VerdierSemanticAnchor{ch3d-1-3-6-morphisms}
\begin{equation}
p_1^i=
\begin{pmatrix}
\operatorname{id}_{X^i}\\0\\0
\end{pmatrix},
\qquad
p_2^i=
\begin{pmatrix}
0\\\operatorname{id}_{X^i}\\0
\end{pmatrix},
\tag{1.3.6.2}
\end{equation}
et il suffit de montrer que $F(p_1)=F(p_2)$. Or le couple de morphismes
$(\operatorname{id}_X,\operatorname{id}_X)$ muni de l'homotopie triviale
définit un morphisme $\pi\colon Z\longrightarrow X$ :

\VerdierSemanticAnchor{ch3d-1-3-6-retraction}
\begin{equation}
\pi^i=(\operatorname{id}_{X^i},\operatorname{id}_{X^i},0)
\tag{1.3.6.3}
\end{equation}
et on a $\pi p_1=\pi p_2=\operatorname{id}_X$. De plus, on vérifie
immédiatement que les morphismes $\pi$, $p_1$ et $p_2$ sont des
homotopismes. Il en résulte que $F(p_1)$, $F(p_2)$, $F(\pi)$ sont des
isomorphismes et que
$F(\pi)F(p_1)=F(\pi)F(p_2)=\operatorname{id}_{F(X)}$. Par suite,
$F(p_1)=F(p_2)$, ce qui achève la démonstration du lemme.

\VerdierSemanticAnchor{ch3d-1-3-7}
\medskip\noindent\textbf{Remarque 1.3.7.}\quad
Il résulte de la proposition (1.3.5) que la catégorie dérivée de
$\mathcal A$ est la solution du problème universel consistant à inverser
les quasi-isomorphismes de $\operatorname{comp}(\mathcal A)$ (chap. II,
2.2.3, b)).
\endgroup
